% Part: first-order-logic
% Chapter: beyond
% Section: intuitionistic-logic

\documentclass[../../../include/open-logic-section]{subfiles}

\begin{document}

\olfileid[fr]{fol}{byd}{il}

\olsection{Logique intuitionniste}

À la différence des logiques du second ordre et d'ordre supérieur,
la logique intuitionniste du premier ordre restreint la version
classique pour représenter un raisonnement davantage « constructif ».
Quelques exemples permettront d'en éclairer les motivations.

Supposons que quelqu'un vienne un jour vous annoncer qu'il a
déterminé un nombre naturel~$x$ tel que, si $x$ est premier,
l'hypothèse de Riemann est vraie, et que, si $x$ est composé, elle
est fausse. Voilà une excellente nouvelle !{} La vérité de
l'hypothèse de Riemann est l'une des grandes questions ouvertes des
mathématiques; cette personne semble avoir ramené le problème à un
calcul, celui qui consiste à déterminer si un nombre donné est
premier ou non.

Quelle est donc cette valeur miraculeuse de~$x$ ? Voici sa
description : $x$ est le nombre naturel égal à $7$ si l'hypothèse de
Riemann est vraie, et à $9$ sinon.

Vous réclamez alors votre argent. Du point de vue classique, cette
description détermine effectivement une unique valeur de~$x$; mais
vous attendiez une valeur de~$x$ donnée \emph{explicitement}.

Prenons un autre exemple, peut-être moins artificiel. Il est
possible d'élever un nombre irrationnel à une puissance rationnelle
et d'obtenir un résultat rationnel : ainsi,
$\sqrt{2}^2 = 2$. Ce qui est moins évident, c'est la possibilité
d'élever un irrationnel à une puissance \emph{irrationnelle} pour
obtenir un rationnel. Le théorème suivant répond par l'affirmative.

\begin{thm}
Il existe des nombres irrationnels $a$ et $b$ tels que $a^b$ soit
rationnel.
\end{thm}

\begin{proof}
Considérons $\sqrt{2}^{\sqrt{2}}$. Si ce nombre est rationnel,
il suffit de prendre $a = b = \sqrt{2}$. Dans le cas contraire, il
est irrationnel. Or
\[
(\sqrt{2}^{\sqrt{2}})^{\sqrt{2}} = \sqrt{2}^{\sqrt{2} \cdot
  \sqrt{2}} = \sqrt{2}^2 = 2,
\]
qui est bien rationnel. Dans ce cas, prenons $a$ égal à
$\sqrt{2}^{\sqrt{2}}$ et $b$ égal à~$\sqrt 2$.
\end{proof}

Est-ce une démonstration valable ? La plupart des mathématiciens
estiment que oui. Il reste pourtant quelque chose d'insatisfaisant :
nous avons prouvé l'existence d'un couple de réels possédant une
certaine propriété sans pouvoir dire \emph{quel} couple la possède.
On peut démontrer le même résultat en exhibant le couple $a$, $b$
dans la preuve : prenons $a = \sqrt{3}$ et $b = \log_3 4$. Alors
\[
a^b = \sqrt{3}^{\log_3 4} = 3^{1/2 \cdot \log_3 4} = (3^{\log_3
  4})^{1/2} = 4^{1/2}= 2,
\]
puisque $3^{\log_3 x} = x$.
De plus, $b$ est irrationnel : si $\log_3 4 = m/n$ pour des
entiers $m,n>0$, alors $3^m = 4^n$, ce que contredit l'unicité de
la décomposition en facteurs premiers.\footnote{Complément éditorial : la source donne le couple explicite sans démontrer l'irrationalité de $b$ ; nous ajoutons cette vérification.}

La logique intuitionniste vise à représenter un raisonnement où
des démarches comme celle de la première preuve ne sont pas admises.
Prouver qu'il existe un $x$ satisfaisant~$!A(x)$ demande de fournir
un~$x$ précis, puis une preuve qu'il satisfait $!A$, comme dans
la seconde démonstration. Prouver que $!A$ ou $!B$ vaut demande une
preuve de l'un ou de l'autre.

Sur le plan formel, on obtient la logique intuitionniste du premier
ordre en restreignant d'une certaine manière un système de
!!{derivation} du premier ordre. Il existe de même des versions
intuitionnistes des logiques du second ordre et d'ordre supérieur.
Du point de vue mathématique, il s'agit simplement de systèmes
déductifs formels; comme nous l'avons dit, ils visent toutefois à
représenter une forme de raisonnement mathématique. On peut y voir
le raisonnement justifié par une certaine conception philosophique
des mathématiques, tel l'intuitionnisme de Brouwer; ou bien un
raisonnement plus « concret » et satisfaisant, dans l'esprit du
constructivisme de Bishop. On peut aussi débattre de la fidélité de
la formalisation à ses motivations informelles. Quelle que soit
notre position philosophique, nous pouvons cependant étudier la
logique intuitionniste en tant que logique formellement définie;
nombre de logiciens mathématiciens y trouvent, pour diverses raisons,
un intérêt.

Une interprétation constructive informelle des connecteurs
intuitionnistes est généralement appelée \emph{interprétation BHK},
d'après Brouwer, Heyting et Kolmogorov. Elle se présente ainsi :
une preuve de $!A \land !B$ consiste en une preuve de $!A$
associée à une preuve de $!B$; une preuve de $!A \lor !B$ consiste
en une preuve de $!A$ ou de $!B$, avec l'indication explicite du
cas qui se présente; une preuve de $!A \lif !B$ est un procédé
qui transforme toute preuve de $!A$ en une preuve de~$!B$; une
preuve de $\lforall[x][!A(x)]$ est un procédé qui fournit, pour
chaque valeur de~$x$, une preuve de $!A(x)$; et une preuve de
$\lexists[x][!A(x)]$ comprend une valeur de~$x$ et une preuve
qu'elle satisfait~$!A$. On peut exprimer cette interprétation
en termes informatiques sous le nom de « correspondance de
Curry--Howard », ou de paradigme « formules comme types » :
on envisage !!a{formula} comme la spécification d'un type de données,
et ses preuves comme des objets calculatoires de ce type, qui
permettent de reconnaître sa vérité.

On présente souvent la logique intuitionniste comme la logique
classique « moins » le principe du tiers exclu. Le théorème
suivant précise cette idée.
\begin{thm}
Du point de vue intuitionniste, les schémas d'axiomes suivants
sont équivalents :
\begin{enumerate}
\item $(\lnot !A \lif \lfalse) \lif !A$.
\item $!A \lor \lnot !A$
\item $\lnot \lnot !A \lif !A$
\end{enumerate}
\end{thm}
Déduire les instances de l'un de ces schémas à partir de toutes les instances de l'un des
deux autres constitue un bon exercice de logique intuitionniste.

Les premiers systèmes déductifs de logique propositionnelle
intuitionniste, proposés pour formaliser l'intuitionnisme de Brouwer,
sont dus, indépendamment, à Kolmogorov, Glivenko et Heyting.
La première formalisation de la logique intuitionniste du premier
ordre, ainsi que de certaines parties des mathématiques
intuitionnistes, est due à Heyting. Si plusieurs schémas valides
classiquement ne le sont pas intuitionnistement, beaucoup d'autres
le sont.

La \emph{traduction par double négation} établit un lien important
entre les logiques classique et intuitionniste. On la définit par
induction comme suit; $!A^N$ est la traduction « intuitionniste »
de la !!{formula} classique~$!A$ :\footnote{Précision éditoriale : pour compléter les clauses affichées, on prend $\lfalse^N\ident\lfalse$ et $(\lnot !A)^N\ident\lnot !A^N$. Un biconditionnel est développé comme la conjonction des deux implications. La source laisse ces conventions implicites.}
\begin{align*}
!A^N & \ident \lnot\lnot !A \quad \text{si $!A$ est une !!{formula} atomique} \\
(!A \land !B)^N & \ident (!A^N \land !B^N) \\
(!A \lor !B)^N & \ident  \lnot\lnot (!A^N \lor !B^N) \\
(!A \lif !B)^N & \ident (!A^N \lif !B^N) \\
(\lforall[x][!A])^N & \ident \lforall[x][!A^N] \\
(\lexists[x][!A])^N & \ident \lnot\lnot\lexists[x][!A^N]
\end{align*}
Kolmogorov et Glivenko avaient établi des versions de cette
traduction pour la logique propositionnelle; pour la logique des
prédicats, elle est due, indépendamment, à G\"odel et à Gentzen.
On a :

\begin{thm}
\begin{enumerate}
\item $!A \liff !A^N$ est démontrable classiquement.
\item Si $!A$ est démontrable classiquement, alors $!A^N$ est
  démontrable intuitionnistement.
\end{enumerate}
\end{thm}

Imaginons maintenant le dialogue suivant. Un mathématicien classique :
« J'ai démontré $!A$ ! » Un mathématicien intuitionniste :
« Votre preuve n'est pas valable. Ce que vous avez réellement
démontré, c'est $!A^N$. » Le premier répond : « Cela me va ! »
À ses yeux, le second coupe les cheveux en quatre, puisque les deux
énoncés sont équivalents. Mais l'intuitionniste insiste sur leur
différence.

La traduction ci-dessus ne concerne que la logique pure; elle ne
résout pas la question des axiomes \emph{non logiques} à retenir
pour les mathématiques classiques et intuitionnistes, ni celle
du rapport entre ces axiomes. Une légère extension du théorème
précédent apporte cependant une information utile :

\begin{thm}
Si $\Gamma$ démontre $!A$ classiquement, $\Gamma^N$ démontre
$!A^N$ intuitionnistement.
\end{thm}

Autrement dit, si $!A$ se démontre classiquement à partir de
certaines hypothèses, $!A^N$ se démontre intuitionnistement à
partir de leurs traductions par double négation.

Pour montrer qu'un énoncé ou !!a{formula} propositionnelle est
valide intuitionnistement, il suffit de donner une preuve. Mais
comment montrer qu'il ne l'est pas ? Il nous faut pour cela une
sémantique correcte et, de préférence, complète. La sémantique
due à Kripke répond à ce besoin.

On peut reprendre la démarche suivie pour la logique classique :
définir la sémantique, puis démontrer la correction et la complétude.
Une différence mérite toutefois d'être signalée. Pour la logique
classique, la sémantique paraissait, en un sens, « évidente »,
comme si elle était contenue dans le sens des connecteurs. La
sémantique de Kripke peut recevoir une motivation intuitive, sans
produire pour autant la même impression d'inévitabilité. En outre,
la notion de !!{structure} classique est mathématiquement naturelle :
on peut considérer les !!{structure}s comme un outil pour étudier
la logique classique du premier ordre, ou cette logique comme un
outil pour étudier les !!{structure}s mathématiques. Les
!!{structure}s de Kripke, en revanche, semblent n'être que des
constructions logiques, sans intérêt mathématique indépendant.

Pour un langage propositionnel, !!a{structure} de Kripke
$\mModel M = \tuple{W, R, V}$ comprend un ensemble~$W$, un ordre
partiel~$R$ sur~$W$ possédant un plus petit !!{element}, ainsi
qu'une assignation « monotone » des variables propositionnelles
aux !!{element}s de~$W$ : si $u\le v$, alors $V(u)\subseteq V(v)$. Intuitivement, les !!{element}s de $W$
représentent des « mondes » ou des « états de connaissance »;
un élément $v \geq u$ représente un « état futur possible »
de~$u$, et les variables propositionnelles assignées à~$u$
correspondent aux propositions connues comme vraies à l'état~$u$.
La relation de forçage $\mSat{M}{!A}[w]$ étend alors cette
interprétation à toutes les !!{formula}s du langage; on lit
$\mSat{M}{!A}[w]$ : « $!A$ est vraie à l'état~$w$ ».
On la définit par induction :
\begin{enumerate}
\item $\mSat{M}{\Obj p_i}[w]$ si et seulement si $\Obj p_i$
  est une des variables propositionnelles assignées à~$w$.
\item $\mSat/{M}{\lfalse}[w]$.
\item $\mSat{M}{(!A \land !B)}[w]$ si et seulement si
  $\mSat{M}{!A}[w]$ et $\mSat{M}{!B}[w]$.
\item $\mSat{M}{(!A \lor !B)}[w]$ si et seulement si
  $\mSat{M}{!A}[w]$ ou $\mSat{M}{!B}[w]$.
\item $\mSat{M}{(!A \lif !B)}[w]$ si et seulement si,
  pour tout $w' \geq w$, $\mSat{M}{!A}[w']$ entraîne
  $\mSat{M}{!B}[w']$.
\end{enumerate}
Un bon exercice consiste à chercher une !!{structure} de Kripke
qui montre que $\lnot (p \land q) \lif
(\lnot p \lor \lnot q)$ n'est pas valide intuitionnistement.

\end{document}
